%% Research Note of the AAS (draft, NOT submitted): EPIC 202126877, eclipses switching on.
%% Source, numbers and provenance: https://github.com/coracle-us/duo-transits (eb_depth_evolution/; see rnaas/numbers.md).
\documentclass{aastex701}

\begin{document}

\title{Eclipses Switching On: Rapid Inclination Change in the Eccentric A-star Binary EPIC 202126877}

\author[orcid=0009-0000-7524-4470]{Jon McBee}
\affiliation{Independent researcher, USA}
\email{jon.mcbee@gmail.com}

\begin{abstract}
EPIC~202126877 (TYC~1883-1721-1) is a 21.775-d eccentric binary of two A stars whose eclipses have
appeared within the last two decades. SuperWASP and KELT saw no eclipse in 2004--2008 ($<4.5\%$, $3\sigma$).
K2 caught a 6.2\% primary and no secondary in 2014. TESS measured a primary 20.6\% deep in 2021, 27.4\% in
2023 and 35.6\% in 2026, and a secondary growing from 4.4\% to 12.5\%. Two independent eclipse models give
an orbital inclination rising from $85.8^\circ$ (2014) to $88.5$--$88.9^\circ$ (2026). The rate rose from
$0.20^\circ$\,yr$^{-1}$ (2014--2021) to $0.31^\circ$\,yr$^{-1}$ (2021--2026). The primary eclipses probably
began in 2007--2011, and could become total, about 40--45\% deep, around 2028--2031. Within single orbits the primary and secondary eclipse times shift in opposite directions by
up to 26\,min, so $e\cos\omega$ varies on month timescales. Both effects point to a close third body that
Gaia DR3 does not detect.
\end{abstract}

\keywords{\uat{Eclipsing binary stars}{444} --- \uat{Trinary stars}{1714} --- \uat{Orbital evolution}{1178}}

\section{Introduction}

Binaries whose eclipses appear or disappear are rare: SS~Lac \citep{Torres2000} and V907~Sco
\citep{Lacy1999} are classic cases. They record nodal precession driven by a third body, and so measure the
architecture of a hierarchical triple. EPIC~202126877 (TIC~86491361; Gaia DR3 3383995253529149056; $G=10.95$)
is an Am star \citep[A5;][]{Shang2022}. It is listed with a single K2 Campaign~0 eclipse and no period by
\citet{LaCourse2015}, and it is detected automatically in TESS by \citet{IJspeert2024} with $P=21.78$\,d.
In a separate note (McBee 2026, submitted) we pointed out that its eclipses are deepening. Here we quantify
the change and its cause. No other study of the system's eclipses exists (ADS full text, SIMBAD, VizieR).

\section{Data and depths}

We used:
\begin{itemize}
\item K2 Campaign~0 (EVEREST and K2SFF, 2014 May; \citealt{Howell2014});
\item TESS Sectors 43--45 (2021) and 71--72 (2023), from TESS-SPOC \citep{Caldwell2020}, QLP \citep{Huang2020},
  and our own TESScut \citep{Brasseur2019} photometry;
\item a 2026 January TESS full-frame segment (BTJD 4055.8--4063.0) with our own photometry;
\item ground photometry: SuperWASP (2004--2008; \citealt{Pollacco2006}), KELT-North (2006--2008;
  \citealt{Pepper2007}) and DASCH plates \citep[1890--1990;][]{Grindlay2009}.
\end{itemize}
ZTF and ATLAS saturate on this star. ASAS-SN V photometry (2014--2018) is too sparse at eclipse phase to
measure depths. The faint neighbor ($G=14.9$ at $21''$) contributes about 3\% of the light in TESS. For
TESS, all depths come from one fixed aperture, with the depth taken as the median flux within $\pm15$\,min of
mid-eclipse.

The primary depth is 6.2\% (2014), $20.6\pm0.1\%$ (2021), $27.4\pm0.2\%$ (2023) and 35.6\% (2026). The
secondary is $<0.03\%$ (2014), $4.4\pm0.1\%$ (2021) and $12.5\pm0.1\%$ (2023); none fell in the 2026 segment.
The secondary-to-primary ratio, which no constant dilution can alter, doubles from 0.21 to 0.46. Folding the
ground data on the TESS ephemeris gives a primary depth below 4.5\% ($3\sigma$, typically $<1.6\%$) in
2004--2008, for any phase offset from $-0.6$ to $+0.3$\,d. DASCH shows no excess of faint points relative to
30 comparison stars, which rules out deep eclipses lasting through most of the century.

\section{Inclination history}

We fit all eleven TESS eclipses jointly with two limb-darkened spheres on a Keplerian orbit.
\begin{itemize}
\item \textbf{Free parameters:} the radii, the surface-brightness ratio and $e\sin\omega$, shared across
  seasons, plus one inclination per season. $e\cos\omega$ is set by each season's secondary phase.
\item \textbf{Fit:} reduced $\chi^2=1.43$, $R_1/a=0.0297$, $R_2/a=0.0391$, $e=0.35$ and $\omega\approx138^\circ$.
\item \textbf{2014:} the K2 primary then fixes the inclination, and the model reproduces the absent secondary.
\end{itemize}
The inclination rises from $85.7^\circ$ (2014.4) to $87.3^\circ$ (2021.8), $88.0^\circ$ (2023.9) and
$88.9^\circ$ (2026.0). The ranges across third light (0--15\%), limb darkening and phase assumptions are
$0.1$--$0.4^\circ$, larger than the formal errors, and they shift all seasons together. A blind re-analysis
with independent code and its own photometry agrees within these ranges: 85.8, 87.2, 87.9 and $88.5^\circ$.
The rate rose from $0.20^\circ$\,yr$^{-1}$ (2014--2021) to $0.31^\circ$\,yr$^{-1}$ (2021--2026). A
constant rate is excluded, but the data do not show whether the rate is still increasing.

Primary eclipses require $i>84.4^\circ$ and secondaries $i>86.5^\circ$. Extrapolating places the primary
onset in about 2007--2011, consistent with the ground limits. A trend that turned around before 2014 would
predict a 4--7\% primary in 2004--2008, which those limits exclude. The secondary onset falls in about
2018--2019. The primary becomes total (flat-bottomed, about 40--45\% deep) at $i\approx89.2$--$89.9^\circ$:
the split between the two radii is poorly constrained. At the current rate that happens around 2028--2031
(Figure~\ref{fig:hist}). With the Gaia distance and the fitted light ratio, the radii
($1.6$ and $2.1\,R_\odot$) and total mass ($\approx4.3\,M_\odot$) are consistent with two A stars.

\section{Eclipse timing and a third body}

The TESS primary times scatter by 10\,min rms about a linear ephemeris ($\chi^2=483$ for 4 degrees of freedom),
consistent across three reductions and two timing methods. The interval from secondary to primary eclipse fell
by $26\pm2.5$\,min within one orbit in 2021 and by about 50\,min within one orbit in 2023, but by only about
3--4\,min per orbit on average between the two seasons. Within each season the
primary and secondary times move in opposite directions, while their mean stays on the period. That pattern is
a changing $e\cos\omega$, not light-travel time, and it varies over months. Tidal and relativistic apsidal
motion are orders of magnitude too slow. Rapid nodal precession together with these dynamical timing variations
therefore requires a close third body; an order-of-magnitude precession estimate suggests an outer period of
$\sim$100--600\,d. Gaia DR3 \citep{Gaia2023} shows a normal RUWE (1.007), no non-single-star solution and no
co-moving wide companion, which mildly favors a low-mass tertiary.

\begin{figure}[ht!]
\plotone{depth_history_summary.png}
\caption{Top: primary (blue) and secondary (green) eclipse depths of EPIC~202126877 from K2 (2014) and TESS
(2021--2026), with SuperWASP/KELT upper limits (2004--2008). Curves show the eclipse model for a linear (solid)
and quadratic (dotted) inclination trend. Bottom: inclination per season, with the thresholds below which
primary and secondary eclipses vanish.\label{fig:hist}}
\end{figure}

\section{Predictions}

On the linear TESS ephemeris ($P=21.774799$\,d, $T_0={\rm BJD_{TDB}}\,2459488.58926$), the next primaries
fall on 2026 Oct~25 22:39, Nov~16 17:12, Dec~8 11:46 and Dec~30 06:21~UTC ($\pm30$\,min, set by the timing
variations).
Each lasts about 7\,h at $V\approx11$, which is within reach of small telescopes. Timings and depths through
2028--2031 will test the predicted totality and whether the rate keeps increasing. Radial velocities and a photodynamical model could then measure
the tertiary's orbit and mass. Data, code and fits are at
\url{https://github.com/coracle-us/duo-transits} (directory \texttt{eb\_depth\_evolution}).

\begin{acknowledgments}
This work uses K2 and TESS data from MAST: K2 light curves
\dataset[doi:10.17909/t9ws3r]{https://doi.org/10.17909/t9ws3r}, TESS-SPOC light curves
\dataset[doi:10.17909/t9-wpz1-8s54]{https://doi.org/10.17909/t9-wpz1-8s54}, QLP light curves
\dataset[doi:10.17909/t9-r086-e880]{https://doi.org/10.17909/t9-r086-e880} and full-frame images
\dataset[doi:10.17909/0cp4-2j79]{https://doi.org/10.17909/0cp4-2j79}. It also uses the SuperWASP public archive,
KELT-North data from the NASA Exoplanet Archive, DASCH, ESA Gaia DR3, VizieR, SIMBAD and ADS.
\end{acknowledgments}

\facilities{Kepler (K2), TESS, SuperWASP, KELT, Gaia}
\software{lightkurve \citep{Lightkurve2018}, astroquery \citep{Ginsburg2019}, batman \citep{Kreidberg2015}, astrocut \citep{Brasseur2019}}

\section*{AI disclosure statement}
The author used an AI assistant, Anthropic's Claude (a large language model), under his supervision. It wrote
and ran the analysis code, including an independent blind re-analysis and adversarial reviews performed by
separate AI agents, and drafted and edited the text of this note. The author reviewed and edited the text,
checked the numbers against the saved analysis outputs, and takes full responsibility for the content. Every
reference was checked against ADS, arXiv, Crossref or VizieR records. The figure was generated by analysis
code from the data, not by a generative AI model.

\begin{thebibliography}{}
\bibitem[Brasseur et al.(2019)]{Brasseur2019} Brasseur, C.~E., Phillip, C., Fleming, S.~W., Mullally, S.~E., \& White, R.~L. 2019, Astrocut: Tools for creating cutouts of TESS images, Astrophysics Source Code Library, ascl:1905.007
\bibitem[Caldwell et al.(2020)]{Caldwell2020} Caldwell, D.~A., Tenenbaum, P., Twicken, J.~D., et al. 2020, RNAAS, 4, 201
\bibitem[Gaia Collaboration et al.(2023)]{Gaia2023} Gaia Collaboration, Vallenari, A., Brown, A.~G.~A., et al. 2023, A\&A, 674, A1
\bibitem[Ginsburg et al.(2019)]{Ginsburg2019} Ginsburg, A., Sip{\H o}cz, B.~M., Brasseur, C.~E., et al. 2019, AJ, 157, 98
\bibitem[Grindlay et al.(2009)]{Grindlay2009} Grindlay, J., Tang, S., Simcoe, R., et al. 2009, ASP Conf. Ser., 410, 101
\bibitem[Howell et al.(2014)]{Howell2014} Howell, S.~B., Sobeck, C., Haas, M., et al. 2014, PASP, 126, 398
\bibitem[Huang et al.(2020)]{Huang2020} Huang, C.~X., Vanderburg, A., P\'al, A., et al. 2020, RNAAS, 4, 204
\bibitem[IJspeert et al.(2024)]{IJspeert2024} IJspeert, L.~W., Tkachenko, A., Johnston, C., \& Aerts, C. 2024, A\&A, 691, A242
\bibitem[Kreidberg(2015)]{Kreidberg2015} Kreidberg, L. 2015, PASP, 127, 1161
\bibitem[LaCourse et al.(2015)]{LaCourse2015} LaCourse, D.~M., Jek, K.~J., Jacobs, T.~L., et al. 2015, MNRAS, 452, 3561
\bibitem[Lacy et al.(1999)]{Lacy1999} Lacy, C.~H.~S., Helt, B.~E., \& Vaz, L.~P.~R. 1999, AJ, 117, 541
\bibitem[Lightkurve Collaboration et al.(2018)]{Lightkurve2018} Lightkurve Collaboration, Cardoso, J.~V.~d.~M., Hedges, C., et al. 2018, Lightkurve: Kepler and TESS time series analysis in Python, Astrophysics Source Code Library, ascl:1812.013
\bibitem[Pepper et al.(2007)]{Pepper2007} Pepper, J., Pogge, R.~W., DePoy, D.~L., et al. 2007, PASP, 119, 923
\bibitem[Pollacco et al.(2006)]{Pollacco2006} Pollacco, D.~L., Skillen, I., Collier Cameron, A., et al. 2006, PASP, 118, 1407
\bibitem[Shang et al.(2022)]{Shang2022} Shang, L.-H., Luo, A.-L., Wang, L., et al. 2022, ApJS, 259, 63
\bibitem[Torres \& Stefanik(2000)]{Torres2000} Torres, G., \& Stefanik, R.~P. 2000, AJ, 119, 1914
\end{thebibliography}

\end{document}
